%\documentclass[12pt,oneside]{amsart}
\documentclass{scrartcl}

%\documentclass{article}

\usepackage{amsthm,amsmath,amssymb}
\usepackage[numbers]{natbib}

%\usepackage[colorlinks]{hyperref}
\usepackage{hyperref}
\hypersetup{
    colorlinks,%
    citecolor=black,%
    filecolor=black,%
    linkcolor=black,%
    urlcolor=black
}
\usepackage{hypernat}
\usepackage[top=2.5cm, bottom=2.5cm, left=2.8cm,
right=2.8cm]{geometry}
\usepackage{graphicx}


%\setlength{\parskip}{0pt}
%\setlength{\parsep}{0pt}
%\setlength{\headsep}{0pt}
%\setlength{\topskip}{0pt}
%\setlength{\topmargin}{0pt}
%\setlength{\topsep}{0pt}
%\setlength{\partopsep}{0pt}

\linespread{1}

%\usepackage{marvosym}

%\textwidth = 6.5 in \textheight = 9.2 in \oddsidemargin = 0.0 in
%\evensidemargin = 0.0 in \topmargin = -0.0 in \headheight = 0.0 in
%\headsep = 0.0in \parskip = 0.0in \parindent = 0.0in
%\def\baselinestretch{0.871}


\newtheorem{theorem}{Theorem}[section]
\newtheorem{proposition}[theorem]{Proposition}
\newtheorem{problem}[theorem]{Problem}
\newtheorem{fact}[theorem]{Fact}
\newtheorem{lemma}[theorem]{Lemma}
\newtheorem{defn}[theorem]{Definition}
\newtheorem{corollary}[theorem]{Corollary}
\newtheorem{remark}{Remark}[section]
\newtheorem{example}[theorem]{Example}

\newcommand{\E}{{\mathbb E}}
\newcommand{\se}{{\mathcal{E}}}
\newcommand{\R}{{\mathbb R}}
\renewcommand{\P}{{\mathbb P}}
\newcommand{\ac}{{\mathcal{A}}}
\newcommand{\A}{{\mathcal{A}}}
\newcommand{\B}{{\mathcal{B}}}
\newcommand{\C}{{\mathcal{C}}}
\newcommand{\M}{{\mathcal{M}}}
\newcommand{\Mf}{{\mathfrak{M}}}
\newcommand{\T}{{\mathcal{T}}}
\newcommand{\Z}{{\mathcal{Z}}}
\newcommand{\K}{{\mathcal{K}}}
\newcommand{\lc}{{\mathcal{L}}}
\newcommand{\Ps}{{\mathcal{P}}}
\newcommand{\G}{{\mathcal{G}}}
\newcommand{\pa}{{\mathring{p}}}
\newcommand{\F}{{\cal F}}
\newcommand{\samp}{{\mathcal{X}}}
\newcommand{\X}{{\mathcal{X}}}
\newcommand{\hi}{{\hat{\Phi}}}
\newcommand{\data}{{\text{data}}}

\newcounter{rcnt}[section]
\renewcommand{\thercnt}{(\roman{rcnt})}

\def\qt#1{\qquad\text{#1}}

\def\argmin{\mathop{\rm argmin}}
\def\argmax{\mathop{\rm argmax}}


\setlength{\parskip}{1.4 \medskipamount}
%\sloppy
%\linespread{1.3}

\begin{document}

\title{STAT 238 - Bayesian Statistics  \\
  Lecture Fourteen} 
\subtitle{Spring 2026, UC Berkeley}

\author{Aditya Guntuboyina}


\date{23 Feb 2026}

\maketitle


We will illustrate Dirichlet-Multinomial inference on a text data
example. This analysis is taken from
\citet{MacKay1995HierarchicalDirichlet}. 

\section{A Text Example}

We observe some text $y_1, \dots, y_n$. The vocabulary consists of the
number of distinct words (including punctuation marks) in the
text. Assume the vocabulary size is $k$.

\subsection{A simple i.i.d model}

The simplest model assumes that $y_1, \dots, y_n$ are i.i.d with an
distribution $P$ that is supported on the vocabulary. The likelihood
then becomes:
\begin{align}\label{lik_unigram}
 \prod_{i=1}^n p_{y_i} = \prod_{i=1}^k p_i^{N_i}
\end{align}
where $N_i$ is the observed count for the $i$-th word in the
vocabulary. We can place a Dirichlet prior on $(p_1, \dots, p_k)$
(e.g., $\text{Dirichlet}(0, \dots, 0)$) which would result in the
$\text{Dirichlet}(N_1, \dots, N_k)$  posterior distribution. One can
then generate new text in the following way:
\begin{enumerate}
\item First generate $(p_1, \dots, p_k)$ from the posterior
  distribution.
\item Then generate new words $y_{n+1}, y_{n+2}, \dots$ i.i.d from the
  multinomial distribution $\text{Mult}(1; p_1, \dots, p_k)$ until a
  specific word (e.g., the period symbol '$\cdot$')) is generated. 
\end{enumerate}
Obviously, this will not result in realistic text. This is because the
i.i.d assumption collapses the text into a bag-of-words
representation, discarding all sequential information. Indeed, in this
model, we are effectively working with the counts 
$N_1, \dots, N_k$ instead of the raw data $y_1, \dots, y_n$. The
likelihood in terms of the counts is:
\begin{align*}
  \frac{n}{N_1! \dots N_k!} p_1^{N_1} \dots p_k^{N_k} \propto
  \prod_{i=1}^k p_i^{N_i}
\end{align*}
which coincides with \eqref{lik_unigram}. As a result, it cannot
capture even the most basic linguistic dependencies (e.g., that
``the'' is likely to be followed by a noun). 


\subsection{AR(1) model}

A slightly better model will be obtained if we replace the i.i.d
assumption with a Markovian one. This is also referred to as the first
order AutoRegressive AR(1) model.

The likelihood here then becomes:
\begin{align*}
  p_{y_1} p_{y_2 \mid y_1} p_{y_3 \mid y_2, y_1} \dots p_{y_n \mid
  y_{n-1}, \dots, y_1}
\end{align*}
We assume now that $p_{y_i \mid y_{i-1}, \dots, y_1} = p_{y_i \mid
  y_{i-1}}$ for each $i = 3, \dots, n$. This gives
\begin{align*}
  \text{Likelihood} = p_{y_1} \prod_{i=2}^n p_{y_i \mid y_{i-1}}. 
\end{align*}
We also assume that $p_{y_1}$ is constant so we don't need to model
it. This gives: 
\begin{align}\label{lik_ar1}
  \text{Likelihood} \propto  \prod_{i=2}^n p_{y_i \mid y_{i-1}}. 
\end{align}
If you don't like the assumption that $p_{y_1}$ is constant, you can
interpret the above as the conditional likelihood of $(y_2, \dots,
y_n)$ given $y_1$.

Just as in the i.i.d model where we could write the likelihood in
terms of counts, we can rewrite the likelihood \eqref{lik_ar1} using
'bi-gram counts'. For each $i$ and $j$ in $\{1, \dots, k\}$, let $x_{j
  \mid i}$ denote the number of times the word $j$ follows word $i$ in
the text. Also let $x_i = \sum_{j=1}^k x_{j \mid i}$ denote the number
of times word $i$ appears as the ``previous word'' in the text.

Then the likelihood \eqref{lik_ar1} is the same as:
\begin{align}\label{lik_ar1_c}
    \prod_{i=1}^k  \prod_{j=1}^k p_{j \mid i}^{x_{j \mid i}}
\end{align}
where $p_{j \mid i}$ denotes the probability that word $j$ appears
as the next word given that the previous word is $i$.

We can also directly write the likelihood in terms of the bigram
counts $(x_{j \mid i}, j = 1, \dots, k)$ and $i = 1, \dots, k$. The
model here is:
\begin{align*}
  (x_{j \mid i}, j = 1, \dots, k) \overset{\text{ind}}{\sim}
  \text{Multinomial}(x_i; p_{j \mid i}, j = 1, \dots, k). 
\end{align*}
so that the likelihood in terms of the bigram counts is: 
\begin{align}\label{lik_bigram}
  \prod_{i=1}^k  \frac{x_i!}{\prod_{j=1}^k x_{j \mid
  i}!} \prod_{j=1}^k p_{j \mid i}^{x_{j \mid i}}. 
\end{align}
The unknown parameters in this model are the probability vectors
$(p_{j \mid i}, 1 \leq j \leq k)$ for each $i = 1, \dots, k$.

\citet{MacKay1995HierarchicalDirichlet} perform Bayesian inference
with the following prior: 
\begin{align*}
  (p_{j \mid i}, j = 1, \dots, k) \overset{\text{i.i.d}}{\sim}
  \text{Dirichlet}(a_1, \dots, a_k)
\end{align*}
for some $a_1, \dots, a_k > 0$. This means that all the $k$
probability vectors $(p_{j \mid i}, 1 \leq j \leq k)$ are assumed to
be independently distributed according to the same Dirichlet
distribution $\text{Dirichlet}(a_1, \dots, a_k)$. The prior density is
therefore
\begin{align*}
  \prod_{i=1}^k \frac{\Gamma(a_1 + \dots + a_k)}{\Gamma(a_1) \dots
  \Gamma(a_k)} \prod_{j=1}^k p_{j \mid i}^{a_j - 1}. 
\end{align*}
The posterior is then:
\begin{align*}
 \prod_{i=1}^k  \frac{x_i!}{\prod_{j=1}^k x_{j \mid
  i}!} \prod_{j=1}^k p_{j \mid i}^{x_{j \mid i}} \times \prod_{i=1}^k
  \frac{\Gamma(a_1 + \dots + a_k)}{\Gamma(a_1) \dots \Gamma(a_k)}
  \prod_{j=1}^k p_{j \mid i}^{a_j - 1} \propto \prod_{i=1}^k \prod_{j=1}^k p_{j \mid i}^{x_{j \mid i} + a_j
  - 1}
\end{align*}
which means that: 
\begin{align*}
  (p_{j \mid i}, j = 1, \dots, k) \overset{\text{ind}}{\sim}
  \text{Dirichlet}(x_{j \mid i} + a_j, j = 1, \dots, k). 
\end{align*}
So the posterior mean estimate of $p_{j \mid i}$ is:
\begin{align*}
  \hat{p}_{j \mid i} = \frac{x_{j \mid i} + a_j}{x_i + a_1 + \dots +
  a_k}. 
\end{align*}
The choice of $a_1, \dots, a_k$ controls the properties of the above
posterior. Instead of plugging in some values for $a_1, \dots, a_k$,
\citet{MacKay1995HierarchicalDirichlet} propose estimating them by maximizing the
marginal likelihood of the data $(x_{j \mid i}, 1 \leq i, j \leq k)$
over all values of $a_1, \dots, a_k$.

The marginal distribution of the data given the Dirichlet
hyperparameters $a_1, \dots, a_k$ is: 
\begin{align*}
&  \P \left\{x_{j \mid i}, 1 \leq i, j \leq k \mid a_1, \dots, a_k
  \right\} \\
&= \int \prod_{i=1}^k  \frac{x_i!}{\prod_{j=1}^k x_{j \mid
  i}!} \prod_{j=1}^k p_{j \mid i}^{x_{j \mid i}} \times \prod_{i=1}^k
  \frac{\Gamma(a_1 + \dots + a_k)}{\Gamma(a_1) \dots \Gamma(a_k)}
  \prod_{j=1}^k p_{j \mid i}^{a_j - 1} d(p_{j \mid i} \text{ all } j,
  i) \\
&=    \prod_{i=1}^k \frac{x_i!}{\prod_{j=1}^k x_{j \mid
  i}!} \frac{\Gamma(a_1 + \dots + a_k)}{\Gamma(a_1) \dots \Gamma(a_k)}
  \int \prod_{j=1}^k p_{j \mid i}^{x_{j \mid i} + a_j - 1} d(p_{j \mid
  i}, 1 \leq j \leq k)\\
&=  \prod_{i=1}^k \frac{x_i!}{\prod_{j=1}^k x_{j \mid
  i}!}\left[\frac{\Gamma(a_1 + \dots + 
  a_k)}{\Gamma(x_i + a_1 + \dots + a_k)} \prod_{j=1}^k
  \frac{\Gamma(x_{j \mid i} + a_j)}{\Gamma(a_j)} \right]. 
\end{align*}
In the machine learning literature, the marginal distribution of the
data given the hyperparameters of the prior is referred to as the
Evidence (see e.g.,
\url{https://en.wikipedia.org/wiki/Marginal_likelihood}). So: 
\begin{align*}
  \text{Evidence}(a_1, \dots, a_k) = \prod_{i=1}^k \frac{x_i!}{\prod_{j=1}^k x_{j \mid
  i}!}\left[\frac{\Gamma(a_1 + \dots + 
  a_k)}{\Gamma(x_i + a_1 + \dots + a_k)} \prod_{j=1}^k
  \frac{\Gamma(x_{j \mid i} + a_j)}{\Gamma(a_j)} \right]. 
\end{align*}
The Evidence can be maximized over $a_1, \dots, a_k$ to get a good
choice of 
the hyperparameters. Note that the factorial terms do not involve $
a_j$s so they can dropped in this maximization:
\begin{align*}
 \text{Evidence}(a_1, \dots, a_k) \propto \prod_{i=1}^k \left[\frac{\Gamma(a_1 + \dots + 
  a_k)}{\Gamma(x_i + a_1 + \dots + a_k)} \prod_{j=1}^k
  \frac{\Gamma(x_{j \mid i} + a_j)}{\Gamma(a_j)} \right]. 
\end{align*}
It is always better to optimize the logarithm of the evidence (rather than the
evidence directly):
\begin{align*}
 & \log \text{Evidence}(a_1, \dots, a_k) \\ &= \text{constant} + \sum_{i=1}^k \left\{ \log \Gamma(A) - \log \Gamma(x_i + A) +
  \sum_{j=1}^k \left[ \log \Gamma(x_{j \mid i} + a_j) - \log
                \Gamma(a_j) \right] \right\}. 
\end{align*}
where $A := a_1 + \dots + a_k$ and $x_i = \sum_{j=1}^k x_{j \mid
  i}$. One can compute the gradient with respect to $a$ 
using the digamma function (see
\url{https://en.wikipedia.org/wiki/Digamma_function}):
\begin{align*}
  \Psi(z) := \frac{d}{dz} \log \Gamma(z). 
\end{align*}
This gives:
\begin{align*}
 \frac{\partial}{\partial a_j} \log \text{Evidence}(a_1, \dots, a_k)
  = \sum_{i=1}^k \left\{\Psi(A) - \Psi(x_i + A) + \Psi(x_{j \mid i} +
  a_j) - \Psi(a_j) \right\}. 
\end{align*}
With these gradients (scipy has an inbuilt function for digamma
calculation), some numerical optimization routine can be used for
maximizing the log-Evidence over $a_1, \dots, a_k$ (see code). 



\begin{thebibliography}{99}

\bibitem[MacKay and Peto(1995)]{MacKay1995HierarchicalDirichlet}
D.~J.~C. MacKay and L.~C. Peto,
\newblock ``A hierarchical Dirichlet language model,''
\newblock \emph{Natural Language Engineering}, vol.~1,
no.~3, pp.~289--308, 1995.

%\bibitem[Rubin(1981)]{Rubin1981BayesianBootstrap}
%D.~B. Rubin,
%\newblock ``The Bayesian bootstrap,''
%\newblock \emph{The Annals of Statistics}, vol.~9,
%no.~1, pp.~130--134, 1981.
  
%\bibitem[Fisher(1934)]{Fisher1934}
%R.~A. Fisher,
%\newblock ``Probability, likelihood and quantity of information in the logic of
%uncertain inference,''
%\newblock \emph{Proceedings of the Royal Society of London. Series~A,
%Containing Papers of a Mathematical and Physical Character}, vol.~146,
%no.~856, pp.~1--8, 1934.

%\bibitem[Efron(2010)]{Efron}
%Efron, B. (2010)
%\textit{Large-scale inference: empirical Bayes methods for estimation,
%testing, and prediction}. Cambridge University Press, 2010.
  
%     \bibitem[Jaynes(2003)]{Jaynes} Jaynes, E. T, (2003)
%  \textit{Probability theory: the logic of science}. Cambridge
%  University Press, 2003.

%  \bibitem[Sinai(1992)]{Sinai} Sinai, Y. G, (1992)
%  \textit{Probability theory: an introductory course}. Springer
%  Verlag, 1992.  

%\bibitem[{deGroot(1986)}]{DeGrootBlackwell}DeGroot, M. H. (1986)
%       A conversation with David Blackwell.
%\newblock \emph{Statistical Science}, \textbf{1}, 40--53.

%         \bibitem[Mosteller(1987)]{Mosteller} Mosteller, F, (1987)
%  \textit{Fifty challenging problems in probability with
%    solutions}. Courier Corporation, 1987.

%    \bibitem[MacKay(2003)]{MacKay} MacKay, D. J. C., (2003)
%  \textit{Information theory, inference, and learning
%  algorithms}. Cambridge University Press, 2003.

%\bibitem[{Lindley(2013)}]{Lindley}Lindley, D. V. (2013)
%   \textit{Understanding Uncertainty}. John Wiley and sons, 2013.


\end{thebibliography}







\end{document}

%%% Local Variables:
%%% mode: latex
%%% TeX-master: t
%%% End:
